\documentclass[11pt,a4paper]{article}

\usepackage[utf8]{inputenc}
\usepackage[T1]{fontenc}
\usepackage[french]{babel}
\usepackage[a4paper,margin=2.4cm]{geometry}
\usepackage{lmodern}
\usepackage{enumitem}
\usepackage{booktabs}
\usepackage{array}
\usepackage[hidelinks]{hyperref}

\newcommand{\code}[1]{\texttt{\small #1}}
\setlength{\parskip}{0.5em}
\setlength{\parindent}{0pt}

\title{\textbf{Guide de soutenance --- Voyage Assistant}\\[0.2cm]
\large SAE2 BUT3}
\author{}
\date{}

\begin{document}

\maketitle
\vspace{-1cm}

\begin{center}
\small Encadrement : Bilal Faye (LIPN, CNRS UMR 7030) --- Promotion 2025--2026\\
Démo : \url{https://optimalako-voyage-assistant.hf.space}
\end{center}

\tableofcontents
\newpage

\section{Principe de la soutenance}

Durée visée : \textbf{15 à 20 minutes} de présentation + questions. L'objectif est que
le \textbf{jury, sans formation en IA}, comprenne ce que fait l'application et pourquoi
nos choix techniques sont pertinents. Trois règles :

\begin{itemize}[nosep]
  \item \textbf{Montrer avant d'expliquer} : on commence par une démonstration concrète.
  \item \textbf{Vulgariser} : chaque terme technique (RAG, embeddings, NER) est expliqué
        par une image simple avant le mot savant.
  \item \textbf{Équilibre} : chaque membre présente \textbf{environ 3 diapositives}, soit
        un temps de parole comparable.
\end{itemize}

\section{Répartition (équilibrée) et déroulé}

\begin{center}
\renewcommand{\arraystretch}{1.3}
\begin{tabular}{>{\bfseries}l p{6.5cm} c}
\toprule
Membre & Rôle dans la soutenance & Diapos \\
\midrule
Dhanoush Kessavane & Introduction + \textbf{vidéo de démonstration} + front-end & 1 à 3 \\
Carl Similien & Architecture, API et base de données & 4 à 6 \\
Noé Cervera & Le « cerveau » : LLM, RAG, embeddings, NER & 7 à 9 \\
Ako Christian & Actions réelles, sources de vérité, déploiement & 10 à 12 \\
\bottomrule
\end{tabular}
\end{center}

Dhanoush ouvre avec une \textbf{vidéo de test} de l'application (réalisée avant la
soutenance), puis présente brièvement le front-end ; on enchaîne ensuite sur les autres
membres. Chacun conclut sa partie en annonçant le suivant.

\section{Plan diapositive par diapositive}

Pour chaque diapositive : le \textbf{titre}, le \textbf{contenu} à afficher (3--5 puces
maximum, peu de texte), et l'\textbf{orateur}.

\subsection*{Partie Dhanoush --- Introduction \& démo}

\textbf{Diapo 1 --- Page de titre.}
Titre « Voyage Assistant », sous-titre « Assistant IA pour le transport », les 4 noms,
l'encadrant, l'URL de la démo. \emph{Dit :} qui nous sommes, le sujet en une phrase.

\textbf{Diapo 2 --- Démonstration (vidéo).}
La vidéo de test : une recherche de trajet, une réservation, puis une question à
l'assistant (ex. « combien de temps dure mon vol ? ») et une modification de billet.
\emph{Dit :} « Voici l'application en conditions réelles » ; commenter la vidéo en direct.

\textbf{Diapo 3 --- L'expérience utilisateur (front-end).}
React, design clair/sombre, \textbf{responsive} (mobile/tablette/ordinateur),
\textbf{multilingue} (FR/EN/ES), \textbf{voix} (dictée + lecture). \emph{Dit :} comment
l'interface a été pensée pour être simple et accessible.

\subsection*{Partie Carl --- Architecture \& données}

\textbf{Diapo 4 --- Architecture full-stack.}
Schéma : navigateur (React) $\rightarrow$ API (FastAPI) $\rightarrow$ base de données ;
couche IA côté serveur. \emph{Dit :} le rôle de chaque brique, le découpage front/back.

\textbf{Diapo 5 --- L'API et la base de données.}
FastAPI, validation Pydantic, ORM SQLAlchemy ; tables \code{users}, \code{trajets},
\code{billets}, \code{reclamations}, \code{chat\_sessions}. Catalogue de
\textbf{plus de 182\,000 liaisons} réelles (avion, train, bateau, bus). \emph{Dit :} on
manipule de vraies données, pas des exemples.

\textbf{Diapo 6 --- Sécurité et robustesse.}
Requêtes SQL paramétrées (anti-injection), échappement HTML (anti-XSS), garde-fous
contre insultes et injections de prompt, vérification d'identité, secrets hors du code.
\emph{Dit :} une application sérieuse doit aussi être sûre.

\subsection*{Partie Noé --- Le cerveau (à vulgariser++)}

\textbf{Diapo 7 --- Pourquoi un LLM ne suffit pas.}
Un LLM (Gemini) écrit très bien mais \textbf{invente} parfois et ne connaît ni nos
données ni l'actualité. \emph{Image :} « un élève brillant mais qui répond de mémoire ».
\emph{Dit :} d'où la nécessité du RAG.

\textbf{Diapo 8 --- RAG et embeddings (le c\oe ur du cours).}
RAG = « examen à livre ouvert » : on \textbf{cherche le bon passage} puis on répond.
Les \textbf{embeddings} transforment le sens d'une phrase en nombres ; deux phrases
proches par le sens ont des nombres proches (« je veux annuler » $\approx$ « conditions
de résiliation »). Outils : FAISS (recherche rapide), puis un \textbf{reranker} qui
re-classe les meilleurs résultats. \emph{Dit :} c'est ce qui rend les réponses fiables.

\textbf{Diapo 9 --- Comprendre l'utilisateur (NER) et qualité.}
Le \textbf{NER} repère automatiquement le nom/prénom dans une phrase libre. On évoque
l'évaluation : on a testé les réponses (factuelles, anti-hallucination, robustesse aux
abus). \emph{Dit :} l'IA est mesurée, pas juste « branchée ».

\subsection*{Partie Christian --- Actions réelles \& mise en ligne}

\textbf{Diapo 10 --- L'assistant agit vraiment.}
Modifier / annuler / réclamer passent par une \textbf{machine à états} déterministe
(intention $\rightarrow$ vérification d'identité $\rightarrow$ action en base
$\rightarrow$ e-mail + billet PDF). \emph{Dit :} le LLM \textbf{rédige}, le code
\textbf{exécute} --- aucune action inventée.

\textbf{Diapo 11 --- Les 3 sources de vérité + déploiement.}
RAG (règles), base de données (billet personnel, calcul de durée), APIs temps réel
(météo, heure, change). Déploiement réel sur \textbf{Hugging Face Spaces} + base
\textbf{Neon}. \emph{Dit :} l'application est en ligne et fonctionne.

\textbf{Diapo 12 --- Bilan et perspectives.}
Toutes les exigences du sujet sont couvertes (tableau) ; bonus : multilingue, voix,
recherche web. Perspectives : plus de documents RAG, notifications, intégration continue.
\emph{Dit :} on conclut et on ouvre les questions.

\section{Antisèche : expliquer simplement (jury non-IA)}

\begin{center}
\renewcommand{\arraystretch}{1.3}
\begin{tabular}{>{\bfseries}p{3cm} p{11cm}}
\toprule
Terme & Phrase à dire \\
\midrule
LLM & « Une IA qui écrit du texte, comme un correspondant très cultivé, mais qui peut se
tromper s'il répond de mémoire. » \\
RAG & « On lui fait passer un examen à livre ouvert : on lui met le bon document sous les
yeux avant qu'il réponde. » \\
Embeddings & « On transforme le sens des phrases en nombres, pour que l'ordinateur
retrouve les idées proches même sans les mêmes mots. » \\
FAISS & « Un moteur qui retrouve très vite les passages les plus proches. » \\
Reranker & « Un second lecteur, plus lent mais plus fin, qui re-classe les meilleurs
résultats. » \\
NER & « Une brique qui repère automatiquement le nom et le prénom dans une phrase. » \\
Machine à états & « Un parcours guidé et sûr : l'IA ne décide pas toute seule d'annuler
un billet, c'est notre programme qui le fait, étape par étape. » \\
\bottomrule
\end{tabular}
\end{center}

\section{Questions probables du jury et réponses}

\begin{itemize}[leftmargin=1.2em]
  \item \textbf{« Comment évitez-vous que l'IA invente ? »} --- On sépare la rédaction
        (le LLM) de l'exécution (notre code) ; et pour les questions, on lui fournit la
        bonne source (RAG, base, API) au lieu de le laisser deviner.
  \item \textbf{« Le RAG, c'est imposé par le cours, qu'avez-vous fait exactement ? »} ---
        Une chaîne vectorielle complète : découpage en \emph{chunks}, embeddings
        \code{gte-small} en local, index FAISS, puis reranking par cross-encoder.
  \item \textbf{« Vos données sont-elles réelles ? »} --- Oui : plus de 182\,000 liaisons
        (avion, train, bateau, bus) ; les billets, comptes et conversations sont stockés
        en base PostgreSQL.
  \item \textbf{« Y a-t-il des données pour le bus et le bateau ? »} --- Oui, tous les
        modes sont présents (l'avion et le train sont les plus fournis ; le bus et le
        bateau ont des catalogues plus petits mais réels et réservables).
  \item \textbf{« Est-ce vraiment en ligne ? »} --- Oui, sur Hugging Face Spaces avec une
        base Neon ; on peut le montrer en direct (prévoir le réveil du Space, $\sim$30~s).
  \item \textbf{« Pourquoi pas un hébergeur classique ? »} --- Le RAG charge deux modèles
        en mémoire ($\sim$1~Go) ; les offres gratuites à 512~Mo plantent, d'où Hugging
        Face (16~Go).
  \item \textbf{« Comment gérez-vous plusieurs langues ? »} --- L'interface est traduite
        (FR/EN/ES) et l'assistant reçoit une consigne de langue pour répondre dans la
        langue choisie.
  \item \textbf{« Sécurité ? »} --- SQL paramétré, échappement HTML, garde-fous contre
        injections de prompt et abus, vérification d'identité, secrets hors du code.
\end{itemize}

\section{Conseils pratiques}
\begin{itemize}[nosep]
  \item Réveiller le Space \textbf{avant} de passer (premier accès $\sim$30~s).
  \item Avoir la \textbf{vidéo de démo en local} (secours si le réseau flanche).
  \item Peu de texte par diapo : on \textbf{parle}, la diapo \textbf{appuie}.
  \item Préparer un billet de test (numéro + identité) pour la démo de modification.
  \item Chronométrer : $\sim$3--4 min par personne.
\end{itemize}

\end{document}
